%% Supervisor-facing thesis document. Chapter 1 is the full introduction from
%% Ha_Vien_Mastersthesis_2026.tex, with cross-references retargeted to chapter
%% labels this file actually has. Chapters 2--8 keep five argument paragraphs
%% each, still one move per paragraph, written as plain prose rather than as a
%% bullet list. Written on the CSM thesis template (HYthesisML.cls).
%% Kept in step with Ha_Vien_Mastersthesis_2026.tex, which uses the same class and options.

\documentclass[english,twoside,openright,a4paper,censored,csm,software-track-2020]{HYthesisML}

\usepackage{csquotes}

% REFERENCES: numeric (Vancouver) style, one of the three styles the template
% offers. sorting=none keeps the numbers in citation order, as IEEEtran did.
\usepackage[style=numeric,bibstyle=numeric,backend=biber,natbib=true,
            sorting=none,sortcites=true,maxbibnames=99,giveninits=true,uniquename=init]{biblatex}
\addbibresource{bibliography.bib}
% List the authors in the bibliography as Lastname F.
\DeclareNameAlias{sortname}{family-given}
% Without a mark of its own the bibliography keeps the last chapter's title in
% the running head of its continuation pages.
\defbibheading{bibliography}[\bibname]{%
  \chapter*{#1}%
  \markboth{#1}{#1}}

% lmodern, textcomp, color, graphicx, epsfig, subfigure, tikz, amsmath, amssymb,
% caption, titlesec and appendix are all loaded by the class itself.
\usepackage[pdftex, plainpages=false]{hyperref} % For hyperlinks and pdf metadata
\usepackage{fancyhdr} % For nicer page headers
\usepackage{longtable} % The front-matter list of acronyms
%\usepackage{wrapfig} % For nice text-wrapping figures (use at own discretion)

\onehalfspacing %line spacing
%\singlespacing
%\doublespacing

%\fussy 
\sloppy % sloppy and fussy commands can be used to avoid overlong text lines

% STEP 2:
% Set up all the information for the title page and the abstract form.
% Replace parameters with your information.
\title{Performance, Safety, and Engineering Trade-offs of Rust and C++ in Nokia RAN Control-Plane Systems}
\author{Ha Vien}
\date{\today}
% Supervisors are printed on the abstract form rather than on the title page.
% TODO before submission: the template asks for the academic title (Prof., Dr.)
% in front of each name.
\supervisors{Prof. Keijo Heljanko, Christoph Luneburg, Petteri Kotimaa}
% The class prints examiners nowhere, so this is recorded but has no output.
% The faculty has not named the second examiner yet.
\examiners{Prof. Keijo Heljanko}
\keywords{Rust, C++, RAN control plane, tail latency, memory safety, fault injection, incremental adoption}
% Where deposited fills itself in from the class. Additional information is the
% study track, taken from the software-track-2020 class option.
\additionalinformation{\translate{\track}}


\classification{\protect{\ \\
\  Software and its engineering $\rightarrow$ Software notations and tools $\rightarrow$ General programming languages\  \\
\  Computer systems organization $\rightarrow$ Real-time systems $\rightarrow$ Real-time system specification\\
\  Networks $\rightarrow$ Network types $\rightarrow$ Mobile networks
}}

% if you want to quote someone special. You can comment this line and there will be nothing on the document.
%\quoting{Bachelor's degrees make pretty good placemats if you get them laminated.}{Jeph Jacques} 


% OPTIONAL STEP: Set up properties and metadata for the pdf file that pdfLaTeX makes.
% But you don't really need to do this unless you want to.
\hypersetup{
    %bookmarks=true,         % show bookmarks bar first?
    unicode=true,           % to show non-Latin characters in Acrobat’s bookmarks
    pdftoolbar=true,        % show Acrobat’s toolbar?
    pdfmenubar=true,        % show Acrobat’s menu?
    pdffitwindow=false,     % window fit to page when opened
    pdfstartview={FitH},    % fits the width of the page to the window
    pdftitle={Performance, Safety, and Engineering Trade-offs of Rust and C++ in Nokia RAN Control-Plane Systems},
    pdfauthor={Ha Vien},
    pdfsubject={Master's thesis, Computer Science},
    pdfcreator={Ha Vien},
    pdfproducer={pdfLaTeX}, % producer of the document
    pdfkeywords={Rust} {C++} {RAN control plane} {tail latency} {memory safety},
    pdfnewwindow=true,      % links in new window
    colorlinks=true,        % false: boxed links; true: colored links
    linkcolor=black,        % color of internal links
    citecolor=black,        % color of links to bibliography
    filecolor=magenta,      % color of file links
    urlcolor=cyan           % color of external links
}

% The same generated macro files Ha_Vien_Mastersthesis_2026.tex reads. The results
% paragraphs quote these rather than carrying their own copies of the numbers, so a
% rerun of bench/refresh_thesis.sh updates this document with the thesis. Macros
% prefixed Bm belong to the second bare-metal host (a different design). Macros
% prefixed Guest belong to the earlier virtualised campaign, retained for comparison.
% Generated by bench/plot_figures.py -- do not edit by hand.
\newcommand{\ReprLoad}{28}
\newcommand{\ReprLoadAch}{28.0}
\newcommand{\ReprLoadShortfall}{0.1}
\newcommand{\SatKnee}{334}
\newcommand{\TrackingLoad}{336}
\newcommand{\NRepeats}{20}
\newcommand{\NLoadLevels}{14}
\newcommand{\NVariants}{4}
\newcommand{\NRuns}{1120}
\newcommand{\PNineNineSpread}{14}
\newcommand{\PFiftySpread}{2.9}
\newcommand{\PFiftyRange}{3}
\newcommand{\PNineNineRange}{460}
\newcommand{\PNineNineMaxCI}{2540}
\newcommand{\MaxOfferedLoad}{784}
\newcommand{\NTests}{10}
\newcommand{\MinAttainableP}{0.0000019}
\newcommand{\CorrectedFloor}{0.000019}
\newcommand{\NSignificant}{2}
\newcommand{\SmallestHolm}{0.00010}
\newcommand{\NExclZero}{4}
\newcommand{\NTailExclZero}{1}
\newcommand{\TailExclZeroRows}{p99 latency under system}
\newcommand{\RobustExclZeroRows}{p50 latency, Peak RSS and CPU utilisation under system}
\newcommand{\PFiftyDiffAbsMax}{1.6}
\newcommand{\PFiftyDiffCIMax}{3}
\newcommand{\RssLangDiffLo}{1.3}
\newcommand{\RssLangDiffHi}{1.7}
\newcommand{\RssLangCILo}{-0.5}
\newcommand{\RssLangCIHi}{1.9}
\newcommand{\CpuDiffAbsMax}{3.4}
\newcommand{\CrossAllocConsistent}{p50 latency, p99 latency, p99.9 latency, Peak RSS and CPU utilisation}
\newcommand{\CrossAllocReversed}{none}
\newcommand{\NCrossAllocReversed}{0}
\newcommand{\RssAllocDiffLo}{55.9}
\newcommand{\RssAllocDiffHi}{56.5}
\newcommand{\BudgetRepeats}{15}
\newcommand{\BudgetCoreCpp}{34.4}
\newcommand{\BudgetCoreRust}{21.3}
\newcommand{\BudgetSvcCpp}{13.52}
\newcommand{\BudgetSvcRust}{12.63}
\newcommand{\BudgetHandoffCpp}{13.49}
\newcommand{\BudgetHandoffRust}{12.61}
\newcommand{\BudgetCoreShareMax}{0.250}
\newcommand{\BudgetCoreFaster}{Rust}
\newcommand{\BudgetCoreRatio}{1.6}
\newcommand{\BudgetHandoffSlower}{C++}
\newcommand{\BudgetBoundVsCore}{101}
\newcommand{\FfiCrossing}{2.54}
\newcommand{\FfiOpRepeats}{15}
\newcommand{\FfiOpNative}{12.3}
\newcommand{\FfiOpSeam}{12.7}
\newcommand{\FfiOpDelta}{+192}
\newcommand{\FfiOpDeltaLo}{-954}
\newcommand{\FfiOpDeltaHi}{+1341}
\newcommand{\FfiOpP}{0.239}
\newcommand{\FfiOpBoundPct}{10.91}

% Generated by bench/complexity.py -- do not edit by hand.
\newcommand{\CppFuncs}{37}
\newcommand{\RustFuncs}{40}
\newcommand{\CppNloc}{511}
\newcommand{\RustNloc}{608}
\newcommand{\CppCcnTotal}{93}
\newcommand{\RustCcnTotal}{79}
\newcommand{\CppCcnMean}{2.5}
\newcommand{\RustCcnMean}{2.0}
\newcommand{\CppCcnMax}{13}
\newcommand{\RustCcnMax}{6}
\newcommand{\CppCcnMaxFn}{\texttt{worker\_loop}}
\newcommand{\RustCcnMaxFn}{\texttt{add\_drb}}
\newcommand{\CcnTotalDrop}{15}
\newcommand{\NlocRise}{19}

% Generated by bench/artefact_facts.py -- do not edit by hand.
\newcommand{\CppUeCtxBytes}{96}
\newcommand{\RustUeCtxBytes}{88}
\newcommand{\UeCtxGapBytes}{8}
\newcommand{\CppUeCtxStatePad}{3}
\newcommand{\CppUeCtxKeyEnd}{50}
\newcommand{\CppUeCtxKeyPad}{6}
\newcommand{\CppRequestBytes}{160}
\newcommand{\RustRequestBytes}{96}
\newcommand{\SignalingMsgBytes}{68}
\newcommand{\SweepUes}{5\,000}
\newcommand{\UeCtxGapTotalKB}{40}
\newcommand{\CppQueueResidentMB}{1.3}
\newcommand{\RustQueueResidentMB}{0.8}
\newcommand{\CppLineCov}{94.4}
\newcommand{\RustLineCov}{97.9}
\newcommand{\RustRegionCov}{99.1}
\newcommand{\ConfHandOps}{223}
\newcommand{\ConfGenOps}{9\,000}
\newcommand{\FfiSeamLines}{231}
\newcommand{\FfiHeaderLines}{67}
\newcommand{\FfiEntryPoints}{11}
\newcommand{\FfiUnsafeFns}{9}
\newcommand{\DeadlockTsanHits}{3}
\newcommand{\DeadlockTsanRuns}{12}
\newcommand{\FaultClasses}{14}
\newcommand{\FaultClassesPrevented}{7}
\newcommand{\FaultClassesBoundary}{7}
\newcommand{\FaultClassesUndetected}{4}
\newcommand{\LostUpdateExpected}{40\,000}
\newcommand{\LostUpdateRustHandled}{28\,714}
\newcommand{\LostUpdateRustLost}{11\,286}
\newcommand{\LostUpdateCppHandled}{27\,320}
\newcommand{\LostUpdateCppLost}{12\,680}
\newcommand{\CppTestBinaries}{6}
\newcommand{\RustTests}{27}

% Generated by bench/adoption_cost.py -- do not edit by hand.
\newcommand{\AdoptRepeats}{15}
\newcommand{\AdoptJobs}{14}
\newcommand{\BuildLibCppWall}{1.9}
\newcommand{\BuildLibRustWall}{1.2}
\newcommand{\BuildLibCppCpu}{2.6}
\newcommand{\BuildLibRustCpu}{3.9}
\newcommand{\BuildLibWallRatio}{1.7}
\newcommand{\BuildLibCpuRatio}{1.5}
\newcommand{\BuildLibWallFaster}{Rust}
\newcommand{\BuildLibCpuCheaper}{C++}
\newcommand{\BuildAllCppWall}{5.7}
\newcommand{\BuildAllRustWall}{8.1}
\newcommand{\BuildAllCppCpu}{35.9}
\newcommand{\BuildAllRustCpu}{64.8}
\newcommand{\BuildAllCpuRatio}{1.8}
\newcommand{\BuildAllFirstPartyCpu}{29.6}
\newcommand{\BuildAllFirstPartyWall}{3.8}
\newcommand{\BuildDepShareCpu}{54}
\newcommand{\BuildAllFirstPartyCpuRatio}{0.8}
\newcommand{\BuildIncCppTu}{1.2}
\newcommand{\BuildIncCppHeader}{5.8}
\newcommand{\BuildIncRust}{3.6}
\newcommand{\BuildIncTuRatio}{3.1}
\newcommand{\BuildIncHeaderRatio}{1.6}
\newcommand{\BuildConfigure}{0.91}
\newcommand{\FootTarget}{conformance\_runner}
\newcommand{\FootCppShipped}{106}
\newcommand{\FootRustShipped}{723}
\newcommand{\FootShippedRatio}{6.8}
\newcommand{\FootCppStatic}{1577}
\newcommand{\FootRustStatic}{723}
\newcommand{\FootStaticRatio}{0.5}
\newcommand{\FootCppDynamic}{106}
\newcommand{\FootRustDynamic}{269}
\newcommand{\FootDynamicRatio}{2.5}
\newcommand{\FootRustStdShare}{63}
\newcommand{\FootCppStdShare}{93}
\newcommand{\FootSharedCpp}{2531}
\newcommand{\FootSharedRust}{5295}
\newcommand{\FootStaticSmaller}{Rust}
\newcommand{\FootDynamicSmaller}{C++}
\newcommand{\FootStaticRatioAbs}{2.2}
\newcommand{\FootCppLocaleShare}{47}
\newcommand{\DepsCore}{0}
\newcommand{\DepsJemalloc}{8}
\newcommand{\DepsHarness}{17}
\newcommand{\DepsUnsafeJemalloc}{623}
\newcommand{\DepsUnsafeHarness}{1064}
\newcommand{\CoreUnsafe}{0}
\newcommand{\DepsJemallocShipping}{5}
\newcommand{\DepsJemallocBuildOnly}{3}
\newcommand{\DepsHarnessShipping}{16}
\newcommand{\DepsHarnessBuildOnly}{1}

% Generated by bench/baremetal_facts.py -- do not edit by hand.
\newcommand{\BmCores}{24}
\newcommand{\BmCounterRepeats}{10}
\newcommand{\BmCpu}{AMD EPYC 7443P}
\newcommand{\BmCpuDiffHi}{7.6}
\newcommand{\BmCpuDiffLo}{0.3}
\newcommand{\BmCpuLower}{Rust}
\newcommand{\BmCpuOneSign}{yes}
\newcommand{\BmCpuSig}{16}
\newcommand{\BmCpuTested}{16}
\newcommand{\BmCycleDelta}{15}
\newcommand{\BmCycleDeltaDir}{fewer}
\newcommand{\BmCycleDeltaP}{0.0117}
\newcommand{\BmDistro}{Ubuntu 24.04.4 LTS}
\newcommand{\BmFarTailCppHi}{228}
\newcommand{\BmFarTailCppLo}{147}
\newcommand{\BmFarTailRustHi}{115}
\newcommand{\BmFarTailRustLo}{88}
\newcommand{\BmFarTailSepHi}{96}
\newcommand{\BmFarTailSepLoad}{32}
\newcommand{\BmFarTailTopLoad}{128}
\newcommand{\BmGovernor}{performance}
\newcommand{\BmGpp}{13.3.0}
\newcommand{\BmInstrDelta}{11}
\newcommand{\BmInstrDeltaDir}{fewer}
\newcommand{\BmInstrDeltaP}{0.002}
\newcommand{\BmIsolcpus}{none}
\newcommand{\BmJeActiveCpp}{1.43}
\newcommand{\BmJeActiveDelta}{-0.64}
\newcommand{\BmJeActiveRust}{0.79}
\newcommand{\BmJeAllocatedCpp}{0.78}
\newcommand{\BmJeAllocatedDelta}{-0.16}
\newcommand{\BmJeAllocatedRust}{0.62}
\newcommand{\BmJeLiveLower}{Rust}
\newcommand{\BmJeMappedCpp}{42.82}
\newcommand{\BmJeMappedDelta}{+2.01}
\newcommand{\BmJeMappedRust}{44.83}
\newcommand{\BmJeMetaDrift}{0.5}
\newcommand{\BmJeMetadataCpp}{4.19}
\newcommand{\BmJeMetadataDelta}{+2.07}
\newcommand{\BmJeMetadataRust}{6.26}
\newcommand{\BmJeOutsideCpp}{4.0}
\newcommand{\BmJeOutsideDelta}{+3.8}
\newcommand{\BmJeOutsideRust}{7.9}
\newcommand{\BmJeResidentCpp}{11.02}
\newcommand{\BmJeResidentDelta}{-1.92}
\newcommand{\BmJeResidentRust}{9.10}
\newcommand{\BmJeRetainedCpp}{27.72}
\newcommand{\BmJeRetainedDelta}{+3.95}
\newcommand{\BmJeRetainedRust}{31.67}
\newcommand{\BmJeRssCpp}{15.08}
\newcommand{\BmJeRssDelta}{+1.74}
\newcommand{\BmJeRssRust}{16.82}
\newcommand{\BmJeRuns}{240}
\newcommand{\BmKernel}{6.17.0-23-generic}
\newcommand{\BmKneeHi}{512}
\newcommand{\BmKneeLo}{448}
\newcommand{\BmLoadHi}{128}
\newcommand{\BmLoadLevels}{8}
\newcommand{\BmLoadLo}{8}
\newcommand{\BmMetastableCv}{37}
\newcommand{\BmMetastableLoad}{512}
\newcommand{\BmMissDelta}{29}
\newcommand{\BmMissDeltaDir}{more}
\newcommand{\BmMissDeltaP}{0.002}
\newcommand{\BmNohzFull}{none}
\newcommand{\BmNumaNodes}{1}
\newcommand{\BmOverloadCv}{0.4}
\newcommand{\BmOverloadFaster}{C++}
\newcommand{\BmOverloadLoad}{640}
\newcommand{\BmOverloadThrCpp}{158}
\newcommand{\BmOverloadThrRatio}{1.64}
\newcommand{\BmOverloadThrRust}{96}
\newcommand{\BmPFiftyDiffHi}{1.1}
\newcommand{\BmPFiftyDiffLo}{0.5}
\newcommand{\BmPFiftyFaster}{Rust}
\newcommand{\BmPFiftyMedian}{74}
\newcommand{\BmPFiftySig}{14}
\newcommand{\BmPFiftyTested}{16}
\newcommand{\BmPNineNineCvMax}{2.3}
\newcommand{\BmPNineNineCvMed}{1.4}
\newcommand{\BmPNineNineMaxCI}{2.8}
\newcommand{\BmPNineNineRange}{3.6}
\newcommand{\BmPinned}{2-17}
\newcommand{\BmPoolSamplesMax}{5\,761\,198}
\newcommand{\BmPoolSamplesMin}{360\,000}
\newcommand{\BmProducers}{8}
\newcommand{\BmRepeats}{15}
\newcommand{\BmRssAllocHi}{3.2}
\newcommand{\BmRssAllocLo}{1.9}
\newcommand{\BmRssLangHi}{1.8}
\newcommand{\BmRssLangLo}{0.7}
\newcommand{\BmRssLargerEffect}{allocator}
\newcommand{\BmRssSig}{16}
\newcommand{\BmRssSpreadMax}{1.7}
\newcommand{\BmRssTested}{16}
\newcommand{\BmRuns}{480}
\newcommand{\BmRustc}{1.98}
\newcommand{\BmShards}{8}
\newcommand{\BmSmt}{off}
\newcommand{\BmSockets}{1}
\newcommand{\BmStealCores}{16}
\newcommand{\BmStealIntervals}{60}
\newcommand{\BmStealMs}{0.0}
\newcommand{\BmStealWindow}{61}
\newcommand{\BmStealWithSteal}{0}
\newcommand{\BmTrackingMax}{448}
\newcommand{\BmVirt}{none}

% Generated by bench/guest_facts.py -- do not edit by hand.
\newcommand{\GuestCores}{16}
\newcommand{\GuestCpu}{Intel(R) Xeon(R) CPU E5-2680 v4}
\newcommand{\GuestDeadlineLoad}{56}
\newcommand{\GuestDistro}{Rocky Linux 8.10 (Green Obsidian)}
\newcommand{\GuestIsolcpus}{none}
\newcommand{\GuestKernel}{7.0.2-1.el8.elrepo.x86_64}
\newcommand{\GuestLoadLevels}{14}
\newcommand{\GuestOverloadThrRatio}{1.8}
\newcommand{\GuestPFifty}{80}
\newcommand{\GuestPNineNine}{4518}
\newcommand{\GuestPNineNineMaxCI}{1480}
\newcommand{\GuestPNineNineRange}{293}
\newcommand{\GuestProducers}{7}
\newcommand{\GuestRepeats}{20}
\newcommand{\GuestRssAllocHi}{34.3}
\newcommand{\GuestRssAllocLo}{32.5}
\newcommand{\GuestRssAllocOverLang}{13}
\newcommand{\GuestRssLangHi}{3.5}
\newcommand{\GuestRssLangLo}{1.8}
\newcommand{\GuestRssLowRunOffset}{31}
\newcommand{\GuestRuns}{1120}
\newcommand{\GuestSatKnee}{336}
\newcommand{\GuestShards}{7}
\newcommand{\GuestSteal}{0}
\newcommand{\GuestThp}{always}
\newcommand{\GuestTrackingLoad}{336}
\newcommand{\GuestVirt}{kvm}


\begin{document}

% Generate title page.
\maketitle

% STEP 3:
% Write your abstract (of course you really do this last).
% You can make several abstract pages (if you want it in different languages)
% by repeating the otherlanguage environment with a different language.

% The abstract form is a fixed-height box, so it is set single-spaced as in the
% template; the 1.5 spacing of the body would cost it a fifth of its lines.
{\singlespacing
\begin{otherlanguage}{english}
\begin{abstract}
      \noindent\textbf{Context.}
      Radio Access Network (RAN) control planes are written mostly in C++ for performance, at a
      well-known cost in memory safety and concurrency. Rust promises comparable performance with
      compile-time safety, but its fit for this latency-critical setting is unproven.
      
      \noindent\textbf{Objective.}
      The thesis compares modern C++ and Rust on a user equipment (UE) Context Management service
      modelled on a real Nokia 5G base-station control-plane component (gNB-CU-CP).
      Average-throughput comparisons leave four dimensions open: tail latency, the C++ defect classes Rust prevents at compile time, the overhead of
      its safety model, and the cost of incremental adoption.
      
      \noindent\textbf{Methods.}
      One language-agnostic specification is frozen first, then implemented twice with identical
      algorithms, concurrency model, allocator, and build settings. A synthetic workload is increased to
      overload, both versions pass the same parity tests, and the evaluation adds fault injection and a
      foreign function interface (FFI) prototype.
      
      \noindent\textbf{Results.}
      Safe Rust rejects all seven injected memory-safety and data-race defect classes at compile
      time. In C++ the same seven build cleanly and show up only at runtime, if a sanitiser catches
      them. Seven further classes fail in both languages. On two of them C++ tooling reports what
      Rust does not. The seam to C++ leaves the core \texttt{unsafe}-free and costs at most
      $\FfiOpBoundPct\%$ of one operation. Performance changes at saturation. Below it the test resolves a few microseconds in Rust's favour, yet an end-to-end call spends only $\BudgetCoreShareMax\%$ on language
      work. Cross-thread handoff takes the rest, so the sweep reflects the design more than the
      language. Timed alone, Rust's operation is $\BudgetCoreRatio\times$ faster. The Rust arms still meet their offered rate to $\TrackingLoadRust$~kops/s, the C++ arms only to $\TrackingLoadCpp$~kops/s. After both have stopped tracking, \OverloadSustainedFaster{} holds
      $\OverloadThrRatio\times$ the throughput of \OverloadSustainedSlower{} in \OverloadThrAgree{}
      interleaved repeats with both allocators. 
      
      \noindent\textbf{Conclusions.}
      The main cost of adopting Rust is engineering effort and the learning curve, not run-time
      performance at nominal load. The gain is freedom from data races, not from concurrency
      defects in general. A nominal-load test would not expose the gap.
\end{abstract}
\end{otherlanguage}
}

% Place ToC. \mytableofcontents is not used here because it switches to arabic
% numbering as its last act, which would number the acronym list 1 and 2 and then
% let \mainmatter restart the body at 1. Front matter is numbered in roman instead,
% as in Ha_Vien_Mastersthesis_2026.tex, and \mainmatter starts the body at arabic 1.
\cleardoublepage
\pagenumbering{roman}
\setcounter{page}{1}
\tableofcontents

% The class already names this list: \nomname is set to "List of acronyms" in
% \captionsenglish. Scoped to the acronyms this document uses, which is a subset of
% the full draft's. No \caption, or the table would be announced in a list of tables
% this document does not have.
\cleardoublepage
\phantomsection
\addcontentsline{toc}{chapter}{\nomname}
\chapter*{\nomname}
\markboth{\nomname}{\nomname}

\noindent Each acronym below is also written out where the text first uses it.

\begingroup
\singlespacing
\footnotesize
\renewcommand{\arraystretch}{1.15}
\begin{longtable}{@{}p{0.15\linewidth} p{0.79\linewidth}@{}}
\hline\hline
Acronym & Expansion \\
\hline
\endfirsthead
\hline\hline
Acronym & Expansion \\
\hline
\endhead
3GPP & 3rd Generation Partnership Project, the body that specifies the 5G standards \\
ABI & application binary interface \\
AI & artificial intelligence \\
AI-RAN & an architecture in which AI and radio access network workloads share one accelerated infrastructure \\
API & application programming interface \\
AWS & Amazon Web Services \\
CISA & Cybersecurity and Infrastructure Security Agency \\
CPU & central processing unit \\
EU & European Union \\
FBI & Federal Bureau of Investigation \\
FFI & foreign function interface \\
FNV-1a & Fowler--Noll--Vo hash function, variant 1a \\
gNB & next-generation NodeB, the 5G base station \\
gNB-CU & gNB central unit \\
gNB-CU-CP & gNB central unit control plane: the system class studied here \\
HDR & high dynamic range, of the latency histogram \\
IMT-2030 & International Mobile Telecommunications-2030, the ITU framework for 6G \\
ITU & International Telecommunication Union \\
NG-RAN & next-generation radio access network \\
PHY & physical layer \\
RAII & Resource Acquisition Is Initialization \\
RAN & radio access network \\
RQ & research question \\
RRC & Radio Resource Control \\
UE & user equipment, an end device such as a phone or a modem \\
UECM & UE context management, the subsystem reimplemented and measured here \\
\hline
\end{longtable}
\endgroup

\cleardoublepage
\mainmatter

% -----------------------------------------------------------------------------------
% STEP 4: Write the thesis.
% Your actual text starts here. You shouldn't mess with the code above the line except
% to change the parameters. Removing the abstract and ToC commands will mess up stuff.
\chapter{Introduction}

Mobile networks are moving from specialised hardware to software on ordinary servers.
Functions that once ran on fixed-purpose, vertically integrated equipment are now
re-implemented as programs on general-purpose machines. The Radio Access Network (RAN),
long the most hardware-bound part of the stack, is part of that shift.
% Checked (make sure everything is understandable)

In 5G, the base station, formally the next-generation NodeB (gNB), is split into a central
unit and one or more distributed units. The central unit is split again into control-plane
and user-plane functions~\cite{3gpp38401,3gpp38300}. Open-RAN pushes the same idea further:
the RAN becomes virtualised, multi-vendor software linked by open
interfaces~\cite{polese2023oran}. The control plane that results, the gNB central unit
control plane (gNB-CU-CP), is a large, stateful, highly concurrent program. It tracks every
connected device, each of which the standards call a user equipment (UE), runs Radio Resource
Control (RRC) signalling~\cite{3gpp38331}, and coordinates mobility and handovers, all within
strict latency limits.

Load at a busy site is high. A single site can run thousands of connection-management and
mobility procedures per second, each of which has to complete inside the millisecond-scale
budget the 5G system sets for control-plane signalling~\cite{3gpp38913}. A deployed control plane may track millions of active
devices at once. As the RAN moves to cloud-native, disaggregated deployments, the
correctness, reliability, and timing of this software increasingly set the quality of the
whole network. An outage or a tail-latency spike here is not a local defect. Every device
the site serves is affected~\cite{dean2013tail,barroso2017microseconds}.

This control-plane software is, today, written mostly in C++. The choice is a
practical one: C++ offers the low-level control, mature tooling, and performance that
latency-critical systems demand. But that performance comes with a well-known cost. C++
places the burden of memory and concurrency management on the programmer, and decades of
experience show that this burden is routinely mishandled, producing use-after-free errors,
double-frees, data races, and related defects that are expensive to find, dangerous in
production, and responsible for a large share of security vulnerabilities in industry
practice~\cite{msrc2019memory,google2022android}. Rust has emerged as the first mainstream
systems language to promise an escape from this trade-off: memory safety and data-race
freedom enforced at compile time on machine-checked
foundations~\cite{klabnik2023rust,jung2018rustbelt}, with performance intended to be
competitive with C++~\cite{balasubramanian2017rust}. Whether that promise holds for a
telecom-grade, latency-critical control plane, where predictability under overload matters
as much as raw speed, is not well understood. This thesis, carried out at Nokia, addresses
the question directly: it develops two functionally equivalent implementations of a
representative RAN subsystem, a UE Context Management service modelled on a real Nokia
gNB-CU-CP component, in modern C++ and in Rust, and compares them under a realistic
synthetic signalling workload.

Three terms recur throughout. The \emph{gNB-CU-CP} is the control-plane software this work is
about. \emph{UE Context Management (UECM)} is the specific subsystem reimplemented and measured.
\emph{Tail latency}, the p99 and p99.9 percentiles rather than the mean, is the property that,
more than average speed, decides whether such a system meets its
deadlines~\cite{dean2013tail,barroso2017microseconds}. Every acronym is written out where the
text first uses it, and the \nomname{} in the front matter collects them.

\section{Motivation and industrial context}\label{sec:motivation}
Three trends make this comparison timely. The first is the cloud-native evolution of the RAN
itself. Moving control-plane functions onto general-purpose servers raises, rather than
lowers, the premium on software quality: the software is longer-lived, evolves continuously,
runs at higher concurrency, and must meet its deadlines without the determinism that
purpose-built hardware once provided. This shift is accelerated by Open-RAN disaggregation,
which turns the base station into virtualised, multi-vendor software connected by open
interfaces~\cite{polese2023oran}, so that the control plane a vendor ships increasingly
competes on software-engineering quality (reliability, predictable timing, security, and the
ease with which a large codebase can be evolved safely) and no longer on purpose-built hardware. In
this setting a latency regression or a memory-corruption defect is not a localised bug but a
network-wide reliability and security risk, and a language-selection question stops being
academic and becomes a concrete engineering decision.

The second trend is that artificial intelligence (AI) is moving from a tool used to tune the
network into a workload the network is built to carry, and the next generation is being
specified that way from the start. International Mobile Telecommunications-2030 (IMT-2030),
the International Telecommunication Union (ITU) framework for 6G, lists artificial
intelligence and communication among its
six usage scenarios and ubiquitous intelligence among its four design
principles~\cite{itu2023imt2030}. The architecture forming around that is AI-RAN, in which
telecommunications and AI workloads are orchestrated together on shared, accelerated
infrastructure, so the sites that run the radio also host distributed edge AI
platforms~\cite{polese2026airan}. The AI-RAN Alliance states the operational requirement
directly: one common infrastructure has to run both kinds of workload concurrently, and their
service-level agreements may differ~\cite{airan2026orchestrator}. That changes what a language
choice has to buy for a control plane. Its software no longer has the host to itself. It
contends for central processing unit (CPU), memory, and cache with inference work whose
arrival pattern it does not
control, and a co-tenant that takes capacity away lowers the load at which the control plane
saturates. Saturation is where the two languages separate most sharply in the measurements of
Chapter~\ref{chapter:eval}. Memory footprint and behaviour past that point are becoming
selection criteria in their own right, and neither is settled by the architecture above the
code.

The third trend is an industry-wide shift toward memory-safe languages. Large
software organisations report that the majority of their serious security vulnerabilities
stem from memory-safety errors, and that adopting memory-safe languages measurably reduces
them~\cite{msrc2019memory,google2022android,vanderstoep2025android,google2024securebydesign}.
The case no longer rests on vendor reporting alone: a group spanning academia and industry has
argued that memory safety now needs standardised terminology and requirements if it is to be
specified and procured at all~\cite{watson2025memorysafety}. This concern has risen to the level of national
policy: cybersecurity agencies now urge software manufacturers to publish ``memory-safe
roadmaps'' and to migrate away from memory-unsafe languages such as C and
C++~\cite{cisa2023memorysafe,oncd2024building}. In the market this subsystem is sold into it has
also become a regulatory matter. The European Union (EU) Cyber Resilience Act places secure-by-design and
lifecycle vulnerability-handling obligations on manufacturers of products with digital elements,
with the vulnerability-reporting obligations applying from September 2026 and most of the
remainder from December 2027~\cite{eu2024cra}. It names no programming language, and it would be wrong to read it as
requiring one. What it changes is that the recurring defect classes of a memory-unsafe language
become a compliance exposure as well as an engineering one. Rust in particular has moved from a niche
language into operating-system kernels, cloud infrastructure, and increasingly into
telecommunications. Production hypervisors such as Amazon Web Services (AWS) Firecracker, which underpins its
serverless platform, are already written in Rust~\cite{agache2020firecracker}. Yet the evidence base for Rust's suitability rests largely on
general-purpose software and on average-case performance benchmarks. The specific question a
telecom vendor must answer, whether Rust can meet the tail-latency and predictability
requirements of a RAN control plane, whether it genuinely prevents the defect classes that
plague the existing C++ code, and what it costs to adopt it incrementally inside a large C++
codebase, has not been answered with direct, controlled evidence. Providing such evidence
for a representative subsystem is the motivation for this work.

\paragraph{Industrial setting.}
This thesis is carried out at Nokia, one of the small number of vendors that build, maintain,
and evolve production RAN software end to end. That market is concentrated in very few hands:
the five largest suppliers took 96\% of worldwide RAN revenue over the first three quarters of
2025, with concentration by Herfindahl--Hirschman Index at a ten-year
high~\cite{delloro2025ran3q,delloro2026ran}. The subsystem it studies, UE Context
Management (UECM), is modelled on a real gNB-CU-CP control-plane component rather than a
textbook example.

The gNB-CU-CP is not a minor component. In the disaggregated architecture of the 3rd
Generation Partnership Project (3GPP) it is the
control-plane half of the central unit, hosting RRC and managing the UE context
lifecycle~\cite{3gpp38401,3gpp38331}, under the millisecond-scale procedure budgets the 5G system
places on control signalling~\cite{3gpp38913} and at the event rates and connected-device counts
a deployed network sustains. Chapter~\ref{chapter:background} places it in the wider architecture.
Because it is written in C++, it inherits both that language's performance and its memory-safety
and concurrency hazards. That is precisely where the three trends above converge, and software of
this class is shipped, evolved, and operated across live operator networks, so the
language-selection decision is a production concern rather than a hypothetical one.

UECM is a deliberate choice of study object within that setting, for the reasons
Chapter~\ref{chapter:context} sets out: it is small enough to reimplement faithfully and reason about in
full, yet it carries the properties that make control-plane code hard. The implementation is
driven entirely by a public, frozen specification rather than by proprietary code, so the
comparison stays reproducible by readers outside Nokia.

\paragraph{A decision with a generational horizon.}
A language choice in software of this class is not revisited often, and it outlives the
generation it was made for. The next one is being specified now: 3GPP has allocated Release 20 to
the 6G studies and Release 21 to the normative work, against an ITU calendar that asks for
technology proposals in early 2029 and the submitted specifications by
mid-2030~\cite{3gpp2024road6g}. That work explicitly covers migration from and interworking with
5G, so a control plane written for the current generation is maintained alongside the next one
rather than retired when it arrives. Nor do the next generation's objectives relax what motivates
this comparison. IMT-2030 names security and resilience as one of its four design principles,
expects the system to be secure by design, defines resilience as the ability to keep operating
correctly during and after a disruptive event, and carries the ultra-reliable low-latency scenario
forward in strengthened form~\cite{itu2023imt2030}. Freedom from whole classes of defect and
predictable behaviour under disturbance are being written into the objectives of the next
generation rather than left to implementers.

The policy dates named above fall inside the same window. The Cybersecurity and Infrastructure
Security Agency (CISA) and the Federal Bureau of Investigation (FBI) classify the
development of new product lines in a memory-unsafe language as a bad practice where the product
serves critical infrastructure, a scope mobile-network infrastructure sits squarely inside, and
they exempt an existing product from the memory-safe roadmap expectation only when its announced
support ends before January 2030~\cite{cisa2025badpractices}. A vendor that defers the language
question until the next generation is specified will not thereby have avoided making it. It will
have made it by default, in the incumbent language, for another decade of deployments.

What this thesis contributes to that decision is bounded, and the boundary is worth stating
plainly. The measurements here are taken on a subsystem modelled on a current-generation
component and evaluated against requirements that are already published. No claim is made about
requirements that have not yet been written. What crosses the generational boundary is not a
latency figure but the language-level comparison itself: Which defect classes a compiler can rule
out, what the safety model costs, and what incremental adoption inside a large C++ codebase
involves. Those questions have the same answers whichever generation the specification above them
names, and they are cheapest to answer before the next codebase is started rather than after.

\paragraph{A concrete failure mode.}
To make the stakes concrete, consider a mobility storm in which many UEs hand over at once. A
single mishandled pointer, a UE context freed on release while an in-flight signalling message
still references it, produces a use-after-free. In C++ this compiles cleanly and may pass
ordinary testing, only to corrupt an unrelated UE's state or crash the worker thread under
precisely the high-concurrency conditions the system exists to handle. The same latent race
can equally show up as an unpredictable tail-latency spike instead of an outright crash. In
safe Rust the same design does not compile: ownership and borrowing reject the dangling
reference at build time. This one scenario captures why the thesis measures both
\emph{safety} (which defects are prevented, and when) and \emph{predictability} (tail
latency under overload) rather than average throughput alone.

\section{Problem statement}
Two properties of this domain make the language choice consequential. First, manual memory and
concurrency management in C++ exposes the software to defects such as use-after-free,
double-free, data races, and iterator invalidation, which are silent at compile time and costly
to find and fix at telecom scale. Second, what governs whether RAN deadlines are met is not
average throughput but worst-case and tail behaviour under load and
overload~\cite{dean2013tail}, a property that is easy to regress and hard to measure correctly.
Rust promises to remove the first class of defect at compile time while claiming performance
competitive on the second, but that claim rests mostly on general-purpose benchmarks rather
than on stateful, concurrent, latency-critical telecom workloads.

Bringing these threads together, the concrete problem this thesis addresses is the following.
Telecom vendors are under growing pressure, both technical and increasingly regulatory, to adopt
memory-safe languages, yet for the RAN control plane they lack \emph{direct, controlled}
evidence on which to base that decision. It is not established whether a Rust implementation of
a stateful, concurrent control-plane subsystem can match modern C++ in \emph{tail latency and
predictability under overload} (as opposed to average throughput), whether Rust's compile-time
guarantees actually eliminate the specific defect classes that affect the existing C++ code,
and what incremental adoption \emph{costs} inside a large C++ codebase. The available evidence
is largely general-purpose and average-case~\cite{balasubramanian2017rust,qin2020rust}, leaving
these three questions (predictability, prevented-defect scope, and migration cost) open for a
representative subsystem. Answering them under controlled, reproducible conditions is the
problem this thesis sets out to solve.

\section{Research questions and hypotheses}\label{sec:rqs}
The thesis is organised around four research questions, ordered to put predictability and
migration ahead of raw average performance:

\begin{description}
  \item[RQ1] How do Rust and modern C++ compare in latency predictability (median, p99,
        p99.9 and p99.99 tail latency, and the maximum) for a RAN control-plane workload
        under increasing load and overload?
  \item[RQ2] Which classes of memory-safety and concurrency defects that occur in C++ are
        eliminated at compile time in the Rust implementation, and at what stage and cost
        are the remaining ones detected?
  \item[RQ3] What is the overhead, if any, of Rust's safety model (bounds checks,
        ownership-driven copies, abstraction cost) on throughput, CPU utilisation, and
        memory behaviour?
  \item[RQ4] What are the practical implications (effort, foreign function interface (FFI)
        interoperability, and risk) of adopting Rust incrementally within an existing C++ RAN
        software ecosystem?
\end{description}

The questions are approached with three prior expectations, stated here so that the evaluation
can confirm or refute them instead of merely describing results after the fact. \textbf{H1:} For
this workload the two languages are performance-equivalent in the operationally relevant region
below saturation, with any difference confined to the far tail. \textbf{H2:} Safe Rust prevents
at compile time the memory-safety and concurrency defect classes that C++ leaves to
best-effort, after-the-fact detection. \textbf{H3:} The run-time cost of C++-to-Rust
interoperability is negligible relative to the cost of the operations invoked, so that the true
cost of incremental adoption is engineering effort rather than performance. H2 is stated so that
it can fail. The defect catalogue behind RQ2 spans classes safe Rust is expected to prevent and
classes it is expected not to, so the study can locate the boundary of the guarantee instead of
only confirming it. These hypotheses are revisited against the evidence in
Chapter~\ref{chapter:discussion}.

\section{Contributions, results, and limitations}
The thesis makes six contributions. Two are artifacts: a reproducible benchmarking framework
for telecom control-plane workloads, and a frozen shared specification with two parity-tested
implementations of a RAN control-plane subsystem, one in C++ and one in Rust. Three are
results. The first is a comparison of latency and resource behaviour that divides at the
saturation knee, where the ordering between the two languages reverses. The second is a
defect-class taxonomy from controlled fault injection that marks where the compile-time
guarantee stops, in both directions. The third is a bound on what a comparison of this shape can
attribute to a language at all, obtained by decomposing an end-to-end operation into the work the
language performs and the cross-thread handoff around it. The last is a set of engineering and
migration guidelines for language selection in cloud-native RAN evolution
(Chapter~\ref{chapter:discussion}). Neither the reversal nor the attribution bound was anticipated.
Both were found while reading results the study had planned to report differently, and
Chapter~\ref{chapter:discussion} records what the first of them cost H1.

The findings divide three ways. On \emph{safety} (RQ2), every injected memory-safety and
data-race defect class is prevented at compile time in safe Rust. The same defects all compile
in C++, and each is caught, if at all, only at run time and only under the right sanitiser, or on
one class under none of them. Seven
boundary classes were chosen deliberately, and neither language prevents them. On two of those
the C++ tooling is ahead, and on one the defect that safe Rust rules out returns as soon as the
program writes \texttt{unsafe}. On \emph{migration} (RQ4), the FFI seam is small and keeps the core
free of \texttt{unsafe}. Measured against the same operation called natively it costs at most
$\FfiOpBoundPct\%$ of one operation, so the real adoption cost is application programming
interface (API) design and not run-time
overhead. Three costs beyond the seam are measured for the same reason, and two of them come back
against the received wisdom: most of the Rust build's extra cost is compiling third-party code and
not the compiler's own work, and the binary-size ordering depends on which linkage policy the two
toolchains are held to rather than on the languages. Both of these results are final.

On \emph{performance} (RQ1, RQ3) the answer divides twice, and drawing both divisions correctly is
where most of the analysis went. The first division separates what the sweep measures from what it
can attribute. An end-to-end operation in this design is $\BudgetCoreShareMax\%$ language work and
a cross-thread handoff takes the rest, so a matched-pair comparison on that path bounds the design
and not the language. The limit holds whichever way such a comparison comes out, and it reaches
any study that reports an end-to-end median without measuring the share. Timed apart from the
handoff, the operation is $\BudgetCoreRatio\times$ faster in Rust, which the sweep cannot see.

The second division is by regime. Below saturation the two implementations are equivalent end to
end in central behaviour, by a margin no deployment would notice. The paired test does resolve a
median difference of $\PFiftyDiffAbsMax\,\mu$s in Rust's favour under both allocators, on an
end-to-end round-trip of $\PFiftyMedian\,\mu$s, so the equivalence is practical and not statistical.
The far tail in the same regime resolves on the reported host and favours \FarTailLower{}. Pooled
over \FarTailSamplesMin{} to \FarTailSamplesMax{} operations per variant, the two languages' p99.99
values do not meet, and the widest gap between variants is $\FarTailRange\,\mu$s. A second
bare-metal campaign at another shard count agrees on the direction
(Chapter~\ref{chapter:eval}). Past saturation the ordering reverses, and where each
implementation stops meeting its offered rate is itself a difference: the Rust arms track to
$\TrackingLoadRust$~kops/s and the C++ arms to $\TrackingLoadCpp$. Deeper still, where all four have
collapsed, the C++ build sustains $\OverloadThrRatio\times$ the throughput of the Rust
build in \OverloadThrAgree{} interleaved repeats under both allocators, and the Rust build spends
$\OverloadCpuPerOpRatio\times$ the CPU on each operation it does complete, which is what rules out
the reading that it is simply doing less. On resources the Rust process holds a few MiB more resident memory, again
under both allocators, in the same direction and the same order of magnitude on all three machines
measured. How it compares
against the cost of the allocator choice is not: the two are comparable on bare metal and the
allocator's was an order of magnitude larger on the virtualised host this study used first, so the
comparison between them is a property of the machine rather than of the allocators.

Four limitations bound all of this: a synthetic workload, one simplified subsystem,
single-author implementation bias, and an allocator footprint whose size moved with the host.
Chapter~\ref{chapter:discussion} treats each in detail. The intended readers are the vendor
making the language decision, the standards discussion around memory-safe
roadmaps~\cite{cisa2023memorysafe,oncd2024building}, and anyone repeating a comparison of this
shape, who can reuse the specification, the pair, the harness and the catalogue.

The scope is drawn deliberately to protect the internal validity of the language comparison. In
scope are the UE Context Management functionality of Chapter~\ref{chapter:context}, the
shared-nothing concurrency design, the load sweep to overload, the fault-injection catalogue,
and the FFI prototype. Explicitly out of scope are the full gNB system and its remaining
control-plane domains, physical layer (PHY) and layer-2 processing and the data-heavy user
plane, real captured traffic
traces (as opposed to a spec-justified synthetic workload), managed-language alternatives
(excluded on the tail-latency grounds argued in Chapter~\ref{chapter:background}), and
production deployment and operational concerns such as upgrade, observability, and
multi-tenancy. These exclusions trade external validity for the controlled conditions a fair
language comparison needs, and each is revisited as a limitation in
Chapter~\ref{chapter:discussion}.

\chapter{Background: Systems Languages for Latency-Critical Software\label{chapter:background}}
% The class prints the chapter titles in the running head in capitals, where the
% full title of this chapter runs into the margin.
\chaptermark{Background}

In the next-generation radio access network (NG-RAN) architecture the base station (gNB) is
split into a central unit and a distributed unit. The central unit is split again into a control plane and a user plane.
UE Context Management sits in the control plane, the gNB-CU-CP, and talks to its
neighbours over the standardised F1, E1, and NG interfaces~\cite{3gpp38401}.

Deadlines are met or missed because of predictability, not average speed. The metrics
that matter are the p99 and p99.9 latencies and the jitter around
them~\cite{dean2013tail}. They must be measured past the saturation point. They must also
be recorded without coordinated omission, that is, without letting a load generator that
has fallen behind hide its own slowest
requests~\cite{schroeder2006openclosed,tene_hdrhistogram}.

C++ offers abstractions that cost nothing at run time, deterministic cleanup through the
Resource Acquisition Is Initialization (RAII) idiom, and no garbage collector. Safety remains the programmer's
responsibility~\cite{szekeres2013sok}. Sanitisers report only the errors that a given
test run happens to reach.

Rust has the same run-time model and also no garbage collector. It checks memory safety
and freedom from data races at compile time, using ownership, borrowing, and thread-safety
traits such as \texttt{Send}~\cite{klabnik2023rust,jung2021safe}. Code that needs to
step outside those rules must be marked \texttt{unsafe}. That keeps such code small and
easy to find.

The comparison is C++ against Rust. Garbage-collected languages are left out because
their pauses are incompatible with the latency targets. The concurrency model is held
fixed in both implementations. The C application binary interface (ABI) is treated as the
realistic adoption path. No
earlier study isolates the language this way for a stateful, concurrent
subsystem~\cite{karlback2021ran,emberg2023asn1}. The five closest studies by domain were
all carried out at Ericsson, and none is peer-reviewed or
indexed~\cite{karlback2021ran,olsson2023ranmgmt,olsson2023asyncpackets,arvidsson2025async,emberg2023asn1}.
Repetition inside one organisation is not generality.

\chapter{Problem Context: UE Context Management and Predictability\label{chapter:context}}
\chaptermark{Problem Context}

UE Context Management owns the lifecycle of one context per connected device. Each
context moves through four states, Idle, Connected, Reconfiguring, and Released, driven
by a small set of lifecycle operations. An operation that is illegal in the current
state returns a typed error. It is not carried out.

The subsystem was chosen because it is representative and still small enough to keep
under control. It has mutable state per device, a strict lifecycle, heavy concurrency,
and hard latency targets~\cite{3gpp38913,3gpp22261}. It is still simple enough to specify
precisely and to test for identical behaviour in two languages.

Seven defect classes fall inside Rust's guarantee. Each has a natural form in this
subsystem: use-after-free, double-free, iterator invalidation, data race, unguarded
shared mutation, use-after-move, and a thread-unsafe handle shared between workers. The
guarantee rests on aliasing control and data-race freedom~\cite{jung2021safe}. The
catalogue is this study's own.

Seven further classes are included on purpose from outside that guarantee: a lock-order
deadlock, a lost update, a leak through a reference cycle, an integer overflow in a
release build, the blast radius of a failed shard worker, a double release across the
C-ABI seam, and a lock guard that is never dropped. They keep the study from confirming a
foregone conclusion. They show where the guarantee stops. Two of them return the opposite
of the expected result. The C++ tooling reports what Rust's does not.

Scale is what makes the tail matter. One procedure triggers many operations. One site
serves many devices. A p99.9 measured on a single operation becomes a steady rate of
degraded procedures in the field~\cite{dean2013tail}. That is why the load sweep is
pushed past the point where the service starts to fall behind.

\chapter{Research Methodology\label{chapter:methodology}}

The approach is empirical and constructive. Instead of comparing two systems that
already exist and differ in many ways, the study builds a matched pair of
implementations from one frozen specification. The language is then the only systematic
difference between them. Both use an identical thread-per-shard concurrency model. That
structuring itself changes tail latency~\cite{enberg2019threadpercore}.

RQ1 and RQ3 are answered with an open-loop load sweep from light load into overload on
pinned cores. Open-loop means the generator keeps to its send schedule instead of
waiting for replies. Latency is recorded in identically configured high-dynamic-range (HDR)
histograms~\cite{tene_hdrhistogram}. Resource use is sampled through the operating
system. The reported sweep runs on a dedicated \MainCpu{}, simultaneous multithreading
\MainSmt{}, frequency governor fixed at \MainGovernor{}, with \texttt{nohz\_full} on the
pinned cores and no hypervisor. That host resolves the overload divergence, the
resource comparison, and the far tail below saturation. The same design was first
measured on a \GuestCpu{} guest (\GuestVirt{}) that could not resolve the far tail:
there the widest interval on one variant's p99 exceeded the spread between all four by
a factor of five, so no ordering could be read however many repeats were run. A third
sweep on a second bare-metal \BmCpu{} uses a different shard count,
load grid, repeat count and toolchain, so it tests whether the findings survive a
change of design rather than a change of machine.

RQ2 is answered by controlled fault injection~\cite{natella2016fault}. A fixed catalogue
of defects is written in both languages. The expected outcome of each one is recorded in
advance. A surprise then shows up as a discrepancy, not as something explained away
afterwards.

RQ4 is answered by a working prototype in which C++ drives the Rust core through a
foreign-function interface. Cross-language memory-management defects concentrate at that
seam~\cite{li2022crosslang}. It yields three measurements: How large the seam is, how
much \texttt{unsafe} code it needs, and what one call across it costs.

Fairness is the central experimental control. The specification is frozen. Allocator and
optimisation settings are matched. Every point where the two languages had to differ is
written down in a forced-idiom log. Reporting follows the guidance that a single figure
invites conclusions run-to-run variation does not
support~\cite{georges2007rigorous}. Medians come with nonparametric bootstrap confidence
intervals~\cite{efron1993bootstrap}. Exact paired permutation tests are corrected for
multiple comparisons by the Holm procedure~\cite{holm1979simple}. All data is synthetic,
so no data-protection concerns arise.

\chapter{Design and Implementation\label{chapter:design}}

Nothing is implemented until the specification is frozen. It fixes the public API, the
UE state machine, the data structures and their algorithmic strategy, the locking and
sharding granularity, and the concurrency model. Both implementations conform to it.
The comparison then measures the language, not a difference of design or of skill. The
specification is itself a deliverable. It is the study's main control against
implementation bias.

One concurrency model is used on both sides, shared-nothing and thread-per-shard.
That structuring is known to change tail latency by
itself~\cite{enberg2019threadpercore}. The identifier space is partitioned into shards.
Each shard is owned by a single worker thread that alone reads or mutates it, so no lock
guards context state. Callers submit a request to the owning shard's bounded queue and
wait on a single-use completion handle. The queue blocks when it is full. That is
deliberate. A drop policy would hide overload, and this study needs overload to appear as
latency. The only shared mutable state either version permits is one atomic counter. No
control flow depends on its value.

Fairness is enforced below the design as well as in it. Both builds link the same
\texttt{jemalloc} and keep a system-allocator build as a control. Both compile at
matched optimisation. Both record latency through one HDR histogram algorithm configured
identically~\cite{tene_hdrhistogram}. Hashing is pinned on both sides, SplitMix64 for
shard assignment and FNV-1a, the Fowler--Noll--Vo hash, inside the map. C++'s default hash of an integer is typically
the identity function. Rust's is cryptographically strong. Left as they came, the two
standard libraries would have been the measurement.

Where the languages cannot satisfy a requirement the same way, the divergence is written
into a forced-idiom log and then measured. Rust takes its bounded blocking queue from
\texttt{sync\_channel}. C++ through C++23 has no such primitive. It assembles one from a
mutex and two condition variables. C++ has no payload-carrying enumeration in this idiom
either. A request becomes a tagged struct holding every argument field any operation might
need, $\CppRequestBytes$~bytes against Rust's $\RustRequestBytes$. The caller answers
through a \texttt{variant} it has to unwrap. The net result is a Rust core
$\NlocRise\%$ longer in code lines whose total cyclomatic complexity is
$\CcnTotalDrop\%$ lower~\cite{lizard}. The branchiest function in either core is the C++
worker loop that dispatches that \texttt{variant}. The Rust library contains no
\texttt{unsafe} at all.

Parity is tested and not assumed. Language-neutral vectors run in both implementations
and their transcripts are diffed operation by operation. A cross-diff cannot tell
``both correct'' from ``both wrong in the same way''. A third implementation in Python
acts as an independent oracle. It has to reproduce all $\ConfHandOps$ hand-derived
expectations before it may generate any of its own. A campaign of roughly $\ConfGenOps$
generated operations found no divergence. The transition matrix is enumerated
exhaustively instead of sampled. The C++ side runs under AddressSanitizer,
ThreadSanitizer and UndefinedBehaviorSanitizer. The Rust side runs under Miri. Line
coverage is reported. A clean sanitiser result is therefore bounded by the code that
actually ran.

\chapter{Experimental Setup and Evaluation\label{chapter:eval}}

The workload is synthetic and spec-justified. A deterministic driver walks each UE
through the connection lifecycle the 3GPP specifications define. Its draw order is pinned,
so the per-operation-kind counts come out identical in both languages for a given seed.
That doubles as a parity signal. Each sweep point is repeated \NRepeats{} times. The
four language and allocator variants are interleaved inside every repeat, so thermal or
background drift falls on all of them alike. Every point is a median with a bootstrap
$95\%$ confidence interval~\cite{georges2007rigorous,efron1993bootstrap}. The
representative operating point is fixed at $\ReprLoad$~kops/s and not derived. The rule
this study began with, the busiest swept load holding p99 under $5$~ms, located the
virtualised host's noise floor and on a quiet host climbs until it collapses onto the
saturation knee, so it is kept as a diagnostic only. The measured throughput plateau is
$\SatKnee$~kops/s, the lowest of the four, reached at an offered
$\TrackingLoad$~kops/s. All of it runs on a dedicated \MainCpu{} pinned to cores
\MainPinned{}. The frequency governor is fixed at \MainGovernor{} with boost disabled.
The kernel runs tickless on those cores. There is no hypervisor. Isolating the cores
with \texttt{isolcpus} is omitted on purpose, because it would collapse the pinned set
onto one core. The same design was first measured on a \GuestCpu{} guest that could not
resolve the far tail. A second bare-metal host at a different shard count confirms that
the directions survive a change of design.

Below saturation the four variants sit within $\PFiftyRange\,\mu$s of one another at the
median, $\PFiftySpread\%$. The paired test resolves a difference of
$\PFiftyDiffAbsMax\,\mu$s in Rust's favour under both allocators. Higher up the
distribution the ordering can now be read. The widest interval on a single variant's own
p99 spans $\PNineNineMaxCI\,\mu$s against $\PNineNineRange\,\mu$s of spread between all
four, so the uncertainty about one is smaller than the gap between them. On the
virtualised host the same two quantities were $\GuestPNineNineMaxCI\,\mu$s against
$\GuestPNineNineRange\,\mu$s, the wrong way round. The far tail resolves here too, from the
repeats' histograms pooled into \FarTailSamplesMin{} to \FarTailSamplesMax{} operations per
variant, and both \FarTailLower{} variants sit below both of the others by up to
$\FarTailRange\,\mu$s. A second bare-metal host at another shard count agrees, running
$\BmFarTailCppLo$--$\BmFarTailCppHi\,\mu$s for the C++ variants against
$\BmFarTailRustLo$--$\BmFarTailRustHi\,\mu$s for the Rust ones, without overlap. A second
limit is structural. No host would fix it. An
end-to-end operation is $\BudgetCoreShareMax\%$ the work the language performs. The rest
is a cross-thread handoff charged to the operating
system~\cite{zhou2025oltp}. This sweep cannot bound a language effect in either
direction. Measured on its own the bare operation is $\BudgetCoreRatio\times$ faster in
\BudgetCoreFaster{}. The likely cause is the standard hash map's layout, not anything
about safety.

Once the offered rate outruns the service, two differences appear. They stop
meeting the offered rate at different loads: the Rust arms hold to
$\TrackingLoadRust$~kops/s and the C++ arms to $\TrackingLoadCpp$, so at
$\KneeGapLoad$~kops/s \KneeGapFaster{} still serves $\KneeGapThrFast$~kops/s where
\KneeGapSlower{} has collapsed to $\KneeGapThrSlow$. That rung is described and not
tested, because the operating point was found after the data was in hand. Deeper, at
$\OverloadLoad$~kops/s offered, the ordering reverses: \OverloadSustainedFaster{} sustains
$\OverloadThrCpp$~kops/s where \OverloadSustainedSlower{} sustains
$\OverloadThrRust$, a factor of $\OverloadThrRatio$, in \OverloadThrAgree{}
repeats under both allocators with the Holm-corrected $p$ at the floor of the whole
test family. The slower variant is not simply doing less work. It draws
$\OverloadCpuPerOpRatio$ times the CPU per completed operation, which is the signature of
effort spent without completing anything. Isolating the bounded queue and the completion
handle on their own does not reproduce that ordering at significance, so the mechanism sits
in the wider handoff path. This is the largest language effect in the tested family. A
comparison that stopped at nominal load would have concluded the opposite.

Safety is measured by injecting $\FaultClasses$ defects. Seven come from inside the
region Rust claims to govern. Seven were chosen because safe Rust is expected \emph{not}
to prevent them. Every one of the first seven is a compile error
in Rust. It cannot be written at all. In C++ each is valid code. The defect appears only
if a test drives it, and no single sanitiser covers them. The use-after-move is the case
where no sanitiser reports anything: a moved-from object is valid but unspecified, so
there is nothing to detect and the identity is simply lost. All seven boundary cases
compile in safe Rust and misbehave. On two of them the C++ tooling is the one that
reports. LeakSanitizer flags the reference cycle that Rust leaks in silence.
UndefinedBehaviorSanitizer flags the overflow that Rust defines to wrap. ThreadSanitizer
catches the lock-order inversion on some runs and not others, at a rate that is itself
unstable. It builds its lock-order graph from acquisitions that succeed. A run whose first
inversion blocks never closes the cycle. Two Rust model checkers do report both original
boundary defects, Loom exploring interleavings soundly and Shuttle sampling the ones Loom
cannot reach~\cite{bornholt2021shardstore}, but only after each was re-imported against
the checker's own primitives and shrunk to model size. Miri reports the seam double release, the one class
that safe Rust rejects and unsafe Rust does not, but it interprets Rust and cannot
execute the C++ host the defect ships behind. That half of the boundary is caught at test
time, and for the concurrency cases only after a rewrite.

Resource use separates the variants where latency could not. Two effects of comparable
size sit in the resident set. The Rust build costs $\RssLangDiffLo$--$\RssLangDiffHi$~MiB
more peak resident set than the C++ build under either allocator. Linking
\texttt{jemalloc} at all costs $\RssAllocDiffLo$--$\RssAllocDiffHi$~MiB in either
language. Both are real. The language effect keeps its direction and order of magnitude on
all three machines measured. The allocator's was an order of magnitude larger on the virtualised host this
study used first, so their ratio is a property of the machine. The language difference is
a fixed cost, not a per-context one, and the quiet host locates it. \texttt{jemalloc}'s
own accounting there puts $\BmJeMetadataDelta$~MiB more in arena bookkeeping and
$\BmJeAllocatedDelta$~MiB \emph{less} in live allocation, and its resident bytes are
$\BmJeResidentDelta$~MiB lower while the process peak is $\BmJeRssDelta$~MiB higher. So
the footprint sits outside the allocator, in the runtime a Rust process carries before it
allocates. Counters at the operation scope explain the timing result instead: the Rust build
retires $\BmInstrDelta\%$ \BmInstrDeltaDir{} instructions in $\BmCycleDelta\%$
\BmCycleDeltaDir{} cycles. CPU inverts with the regime. It is marginally lower
for Rust while the service copes and much higher per completed operation once it does
not. The migration seam is small on every axis measured: $\FfiEntryPoints$ C entry
points in $\FfiSeamLines$ lines of Rust. All raw-pointer work is confined to
$\FfiUnsafeFns$ functions. The core is left free of \texttt{unsafe}. No cost is detected
against the same operation called natively. That bounds any per-operation tax at
$\FfiOpBoundPct\%$. Three adoption costs beyond the seam were measured as well. Two of
them contradict the received objection. $\BuildDepShareCpu\%$ of the Rust
clean build is third-party code. The project's own source compiles in
$\BuildAllFirstPartyCpuRatio$ times the CPU \texttt{g++} needs. The
$\FootShippedRatio\times$ binary-size gap is a linkage policy that reverses when the
policy is held fixed. The compared component resolves $\DepsCore$ crates and contains
$\CoreUnsafe$ \texttt{unsafe} sites.

\chapter{Discussion\label{chapter:discussion}}

What a deployment should do with RQ1 depends on where it expects to sit, not on a single
ranking. A service sized to run below every arm's knee should prefer the arm that holds
its offered rate further. One that a signalling storm will drive past every knee should
prefer the arm that still completes more work afterwards. Those two preferences point at
opposite languages, which is why Chapter~\ref{chapter:eval} reports both halves rather
than a winner. The closest prior work measures this quantity for 5G control-plane
procedures. It finds the same shape and traces it to stackless
\texttt{async}/\texttt{await}~\cite{arvidsson2025async}. That attribution is accepted
here. The shape appeared with no asynchronous runtime on either side and the concurrency
model held fixed. A runtime cannot be the only place such an effect originates. The
prediction registered in advance also named where to look first, bounds and overflow
checking. That part of it was wrong. Recording the miss is worth more to the next study
than quietly substituting the explanation that survived.

On safety the difference is one of kind. The boundary is reported as a first-class
result, not as a caveat. A compiler that refuses to build a program rules the defect out
of every execution. A sanitiser rules it out only of the executions an instrumented test
happens to take. What is absolute there is the execution coverage, not the breadth of
the catalogue. Rust reasons about aliasing and not about acquisition order. A
\texttt{Mutex} guard whose lifetime the borrow checker tracks perfectly says nothing
about whether two threads take their locks in a compatible
order~\cite{zhang2026deadlock}. The lost update escapes ThreadSanitizer for the same
reason it escapes the borrow checker. Where every access is correctly synchronised, there
is no unsynchronised access to find. Adopting Rust retires one category of concurrency
defect and leaves untouched the category that fails silently. A migration plan reading
``memory-safe'' as ``concurrency-safe'' would mis-budget its review effort. One case
reaches past the injected programs and narrows the credit due on both sides. A
specification gap let an operation be submitted outside the service's running window. It
survived every test, every parity vector and all three sanitisers because nothing ever
called that path. That is the weakness of detection conditional on coverage. What turned
undefined behaviour into a diagnosable failure on the Rust side was an ordinary bounds
check, not a proof. The property worth budgeting for is ``checked by default'', not
``ruled out at compile time''. The root cause belongs to the method. A prohibition stated
without a mechanism behind it leaves each implementation to fail in its own idiom.

RQ3 splits by resource and the splits do not point the same way. RQ4 comes out cheap in
the machine and expensive in the people. The memory penalty reproduces under both
allocators and on a second machine with a different compiler. It is credible beyond its
$p$-value. Peak resident set is a run-level maximum insensitive to the jitter that
destroys a p99. It is a few MiB of fixed cost, and sized against the allocator's own
effect the two are of the same order on both bare-metal hosts, interleaved on the reported
one and with the allocator ahead on the other. On the virtualised host the allocator's was
an order of magnitude larger. The language effect travels and the ratio
does not. The quiet host also locates
the cost. \texttt{jemalloc}'s own accounting puts it \emph{outside} the allocator,
because the Rust build holds less live data and fewer resident bytes there and still
peaks $\BmJeRssDelta$~MiB higher. What a Rust process carries before it allocates is the
standard library and the runtime. Which of its pages hold the difference is not
attributed to a call site. One close study reports the
opposite sign, with Rust lowest and the gap widening as concurrency
rises~\cite{karlback2021ran}. The two are compatible. Their C++ runs on Boost.Fibers,
which allocates a stack per fiber. Their Rust runs on Tokio, where a task is a small
heap-allocated future. What they measured is the per-context cost of two concurrency
runtimes. Two further studies agree with the direction found
here~\cite{olsson2023ranmgmt,sundvisson2026drivers}. On migration the dominant cost is
human~\cite{fulton2021adopting,crichton2024learning}. Confining \texttt{unsafe} to the
seam bounds what has to be reviewed, not what can be corrupted. Host and component share
one address space~\cite{rivera2021galeed,kirth2022pkrusafe}. The build and footprint axes
point the same way as the seam. An organisation worried about either should audit its
dependency budget before blaming the compiler. The one axis where C++ leads outright is
the edit-compile loop for changes that stay inside a translation unit.

The three hypotheses were committed before any measurement. Stating their disposition is
what separates a test from a description. H1's median clause holds at the representative load. Its far-tail clause does not, and
overload, which the wording never named, is where the languages actually part. The
reported campaign now resolves that tail, in Rust's favour, and a second bare-metal host at
another shard count agrees on the direction. H1 was right about where a below-saturation
difference would show up and wrong that it was the only place. It was also not well posed.
``Performance-equivalent'' reads as one quantity and names at least three. H2 holds
without exception for the seven injected classes of its kind. It fails in the looser
reading its wording invites. It is accepted for memory safety and data-race freedom,
rejected for concurrency correctness at large, and it is a claim about \emph{safe} Rust
only: the seam case shows the same class returning inside \texttt{unsafe}. H3 is supported by the bound on a cost the
measurement could have missed, not by the failure to detect one. Three threats stay
load-bearing behind all of it. One author wrote both implementations. The parity suite
bounds behaviour and not performance. The end-to-end instrument
measures the design faithfully and the language barely at all. Single-author bias is
documented and not hypothetical: the closest matched pair in this industry reversed
its own verdict across two benchmarking sessions, comparable after a code review and then
Rust worse by $85$--$91\%$ after further optimisation~\cite{emberg2023asn1}. The overload divergence
is replicated here with its mechanism narrowed and not closed, and no peer-reviewed study
reports the same reversal.

The guidelines are the practical form of all this. Each carries the limit of the
finding behind it. Choose a first component for the narrowness and the scalarity of its
interface, not for its size. The redesign of a rich interface is the schedule risk. The
boundary itself is not. Keep the C ABI in its own crate and the core free of
\texttt{unsafe}. That keeps unsafety out of the dependency chains it otherwise propagates
through~\cite{evans2020rust}, and makes the reviewable surface a countable thing. Leave
the concurrency review budget where it is. Prefer a
shared-nothing structure in either language. The strongest concurrency result here
belongs to the design and to neither compiler. Plan the learning curve as a training
cost and instrument it. The second component can then be scheduled on measured effort.
Measure the component as well as the round-trip. Sweep past the point where the service
stops coping. That is the guideline this study would most like to have followed from the
start.

\chapter{Conclusion and Future Work\label{chapter:conclusions}}

The thesis compared modern C++ and Rust on a UE Context Management service modelled on a
real gNB-CU-CP component. It looked at four dimensions: tail-latency predictability,
compile-time defect prevention, safety-model overhead, and incremental-adoption cost.
Both implementations were built against one frozen specification and driven along identical
code paths. That is what lets a language effect be told apart from a difference of design
or of skill.

On safety the difference is one of kind and it has a clear boundary. The edge stops where
the catalogue says it does. Safe Rust rejects all seven injected memory-safety and
data-race classes at compile time. Hazards that C++ can only hope to catch with
instrumented tests become programs that cannot be written. All seven boundary classes
compile and misbehave in safe Rust as they do in C++. On two of them the C++ tooling
reports what Rust's does not. Stated precisely, Rust eliminates data races and aliasing
errors. It does not eliminate race conditions, lock-ordering errors, leaks, wrapping
arithmetic or the blast radius of a failed worker, and it eliminates nothing inside an
\texttt{unsafe} block.

On performance the evidence first divides by what is being measured. An end-to-end
operation here is $\BudgetCoreShareMax\%$ language. The rest is an operating-system
handoff. An end-to-end null bounds nothing about either compiler. The same operation
measured apart is $\BudgetCoreRatio\times$ faster in \BudgetCoreFaster{}. At the
representative load the paired test resolves $\PFiftyDiffAbsMax\,\mu$s in Rust's favour,
too small for a deployment to care about at the median. The far tail on the reported host
orders the languages in the same direction. A virtualised host could not settle that
percentile. The Rust arms then hold their offered rate to $\TrackingLoadRust$~kops/s against
$\TrackingLoadCpp$, and once all four have collapsed \OverloadSustainedFaster{} sustains
$\OverloadThrRatio\times$ the throughput of \OverloadSustainedSlower{}, which spends
$\OverloadCpuPerOpRatio\times$ the CPU on each operation it completes. On memory the Rust
build carries a few MiB more fixed footprint under both allocators, in the same direction
and order of magnitude on all three machines measured. The allocator choice costs about as
much on both bare-metal hosts, and an order of magnitude more on the virtualised one this
study used first.

What follows is a narrow claim, narrower than this thesis expected to make. An unmodified
Rust core can be driven from C++ across a thin seam. The cost on a real operation is
bounded at $\FfiOpBoundPct\%$. All \texttt{unsafe} is confined to the seam. The core is
verified free of it. For stateful, concurrent, partitionable control-plane services of
this kind, at the loads where the service still meets its offered rate, the choice is
which classes of defect an organisation wants to make unwriteable, and whether it will
pay the engineering and learning cost of getting them. Past that point it is a
performance decision as well. What the thesis adds is an instrument, not a verdict: a
frozen specification, a parity-tested pair of implementations, a measurement framework
that reaches past the saturation knee, a decomposition that says how much of what it
measures the language is answerable for, and a defect catalogue that maps where the
guarantee stops.

Future work answers the limitations one at a time. The immediate priority needs no new
machine: vary the transparent-hugepage policy on one host with everything else held
fixed. All three hosts agree that the Rust build carries the larger footprint. They
disagree about how that compares against the cost of the allocator, which the virtualised
host put an order of magnitude higher and neither bare-metal host does. The hugepage
policy is the candidate this study can name but not test, because it differed between the
hosts and only the system-allocator arm inherits it. Beyond that, replaying
anonymised production traces would validate the operation mix. A subsystem with
genuinely shared mutable state would exercise the two languages differently than a
shared-nothing design does. Evaluating Rust async runtimes against the thread-per-shard
model would separate the runtime from the language in the tail
result~\cite{arvidsson2025async}. Measuring what in-process isolation of the two heaps
costs would price what protecting an \texttt{unsafe}-free core is
worth~\cite{rivera2021galeed}. A longer case study tracking defect density and review
effort during an incremental adoption is what would settle the human cost this study can
only name.

% Declaration of generative-AI use, required by the University when a language model has
% been used. Shared with Ha_Vien_Mastersthesis_2026.tex: edit ai_statement.tex, not a copy.
%% Declaration of generative-AI use. The University of Helsinki requires a student who
%% has used a language model in preparing a thesis to state in writing which model was
%% used and in what way; the model may not be named as an author or cited as a source.
%%
%% Shared verbatim by Thesis_POC.tex and Thesis_v1.tex so the two cannot drift apart.
%% It uses no generated macros and no chapter numbers, because the two documents number
%% their chapters differently.
%%
%% The order of the paragraphs follows the University's basic reporting instruction:
%% what the use covered, what it was for, which tool, how it performed and how its
%% output was checked, and what the author took from it.

\cleardoublepage
\phantomsection
\addcontentsline{toc}{chapter}{Statement on the Use of AI Tools}
\chapter*{Statement on the Use of AI Tools}
\markboth{Statement on the Use of AI Tools}{Statement on the Use of AI Tools}

\noindent\textbf{Tool.} Anthropic's Claude Opus~5, used through the Cursor development
environment, is the only generative AI system used in the preparation of this thesis. It was
used throughout the work, from June to August~2026. Following University of Helsinki policy the
model is neither named as an author nor cited as a source: it cannot take responsibility for
what it produces, and it cannot reliably vouch for the literature it paraphrases.

\noindent\textbf{What the assistance covered.} It was extensive, and is declared as such. The
model drafted and revised the running prose of this thesis and proposed the order in which its
arguments are made. It wrote and refactored the two implementations of the UE Context
Management service, the load generator, the fault-injection catalogue, and the interoperability
prototype. It wrote the analysis code behind the bootstrap confidence intervals, the exact
paired permutation tests, and the Holm correction, together with the figure-generating and
build scripts and the \LaTeX{} sources. It also suggested candidate literature.
% ==================== ACTION REQUIRED BEFORE SUBMISSION ==========================
% The abstract of Thesis_POC.tex was drafted with model assistance. The University's
% thesis guide states that a language model may not be used when preparing the thesis
% summary, and in this programme the summary also serves as the maturity test, where
% AI use is never permitted. Rewrite the abstract unaided before submitting; until
% that is done the closing sentence of this paragraph is not yet true of Thesis_POC.
% =================================================================================
The abstract was written by the author unaided, as the University's thesis guide requires of
the thesis summary.

\noindent\textbf{What the model did not settle.} The research problem, the four research
questions, and the three hypotheses are the author's, as is the control on which the study's
value rests: freezing one language-neutral specification and one concurrency model, so that the
language remains the only systematic difference between the two implementations. So is the
decision that keeps RQ2 falsifiable, namely to carry two defect classes -- a lock-order deadlock
and a lost update -- that lie outside Rust's compile-time guarantee and can therefore embarrass
the premise rather than confirm it. Every measurement reported here was run by the author on the
hosts documented in this thesis; no measured value came from a language model. Each measured
number reaches the text through a macro regenerated from the raw logs, so a figure cannot be
transcribed by hand, by the author or by the model. The readings of those numbers are the
author's as well, including the two restraints the evidence imposes: that a far tail a
virtualised host cannot resolve must not be reported as equivalence, and that a sweep dominated
by cross-thread handoff cannot bound an effect of the language.

\noindent\textbf{How the output was checked.} Model output was treated as a draft to be
falsified rather than as a result. Claims were checked against the sources cited for them: every
bibliography entry was resolved against Crossref, DataCite, or the issuing body, and each was
traced to the page or clause carrying the claim made of it. That audit found four passages where
the prose overstated or misplaced its source, one of them a sentence inside quotation marks that
does not appear in the work it was attributed to; all four were corrected, and the audit is kept
with the artefacts so that the correction remains reviewable. Code was checked by running it:
the two implementations are held to one conformance suite whose transcripts are diffed line by
line, and the fault catalogue fixes the expected outcome of each defect before it is injected,
so that a surprise appears as a discrepancy instead of being explained away afterwards. Nothing
is asserted here on the model's authority.

\noindent\textbf{What the author took from it.} The model is dependable at mechanical work --
transposing an idiom between two languages, restructuring a table, keeping a build honest -- and
undependable in a specific and predictable way: it will supply a plausible citation for a claim,
a plausible cause for a measurement, and a confident sentence where the data supports a hedge.
Its fluency is not evidence. The lasting lesson of the work was accordingly less about either
language than about verification. The effort saved in drafting is spent again, and rightly, in
checking, and the parts of this thesis that are defensible are the parts that were checked
against a source, a specification, or a measurement.

\noindent\textbf{Responsibility.} The author is responsible for the whole of this thesis: its
language, its facts, its code, its measurements, and its conclusions. Where it is wrong, the
error is the author's.


% STEP 5:
% The bibliography style is selected with the biblatex options in the preamble.

\cleardoublepage %fixes the position of bibliography in bookmarks
\phantomsection

\addcontentsline{toc}{chapter}{\bibname} % This lines adds the bibliography to the ToC
\printbibliography

\end{document}
